%!TEX program = lualatex
%!TEX encoding = UTF-8 Unicode  
%!TEX root = chapter08.tex  
\documentclass[main.tex]{subfiles}
\begin{document} 
%----------------------------------------------------------------------------------------
%	 
%----------------------------------------------------------------------------------------
\cleardoublepage
\loadgeometry{marginlayout}  
\pagestyle{scrheadings} % Print headers again  
\KOMAoptions{
	headwidth=textwithmarginpar,
	headsepline=true,
	headsepline= 0.5pt,
	footwidth=textwithmarginpar,
}   
\onehalfspacing 
\setcounter{chapter}{7}


\chapter{Calculs avec les radicaux}


\section{La racine carrée} 
\begin{example}\hfill 
\begin{enumerate}[label=\arabic*),leftmargin=*]
	\item Deux nombres ont pour carré $49$ : $(-7)^2=(7)^2=49$. 
	
On dit que $7=\sqrt{49}$ est la racine carré de $49$.
\item Le nombre $-49$ est le carré d'aucun nombre. La racine carrée de $-49$ n'exite pas : l'écriture $\sqrt{-49}$ n'a pas de sens. 
\item Deux nombres ont pour carré $50$ : la racine carré de $50$ et son opposé :
\setlength{\abovedisplayskip}{3pt}
\setlength{\belowdisplayskip}{3pt} 
\[	
(\sqrt{50})^2=(-\sqrt{50})^2=50
\] 
\item Deux nombres ont pour carré $10$ : la racine carré de $10$ et son opposé :
\setlength{\abovedisplayskip}{3pt}
\setlength{\belowdisplayskip}{3pt} 
\[	
(\sqrt{10})^2=(-\sqrt{10})^2=10
\] 
\item $0$ est le carré de $0$ : la racine carrée de $0$ est égale à $0$.
%\item  $\sqrt{36}=6 $, $49^{\numprint{0.5}}=7$
%\item  $\sqrt{-36}$ et $(-49)^{\numprint{0.5}}$ n'est pas défini.
%\item $-\sqrt{36} = -6$ et $-49^{\numprint{0.5}}=-7$. 
\end{enumerate} 
	\end{example}

\begin{definition} Pour tout réel $a>0$. Il existe deux nombres réels dont le carré vaut $a$ :
	\marginnote{
		\begin{tBox}[leftline=false]
			à retenir : pour tout $a>0$ 
			\setlength{\abovedisplayskip}{3pt}
			\setlength{\belowdisplayskip}{3pt} 
			\[	
			\left(\sqrt{a}\right)^2=\left(-\sqrt{a} \right)^2 = a 
			\]   
		\end{tBox}
	}[0cm]
	\begin{enumerate}[leftmargin=*,label=\textbullet]
		\item le nombre \textbf{positif} noté $\sqrt{a}=a^{\frac{1}{2}}$, c'est \og la racine carrée de $a$\fg.
		\item et son opposé $-\sqrt{a}=-a^{\frac{1}{2}}$
	\end{enumerate}
%	La \og racine carré de $a$ \fg{} est l'\textbf{unique réel positif ou nul}  dont le carré vaut $a$. On le note $\sqrt{a}=a^{\frac{1}{2}}$.
%	\setlength{\abovedisplayskip}{3pt}
%	\setlength{\belowdisplayskip}{3pt} 
%	\[	
%	\left(\sqrt{a}\right)^2=\left( a^{\frac{1}{2}} \right)^2 = a 
%	\] 
\end{definition}  
\begin{theorem} Pour tout $n\in\mathbb{N}$, sa racine carrée $\sqrt{n}$ est égale à un entier  ou un irrationnel.
\end{theorem}
\begin{example}\hfill 
	\begin{enumerate}[label=\alph*),leftmargin=*]
	 \item $\sqrt{36}=6$, $\sqrt{225}=15$...
	 \item $\sqrt{2}\in \mathbb{R}\cap \overline{\mathbb{Q}}$,  $\sqrt{5}\in \mathbb{R}\cap \overline{\mathbb{Q}}$ ...
	\end{enumerate} 
\end{example} 
\begin{remark} Objectif du chapitre est d'apprendre à simplifier les écritures de sommes, de produits et de quotients d'expression de la forme $a+b\sqrt{c}$, avec $a$, $b$ et $c$ des entiers. 
\end{remark}
%----------------------------------------------------------------------------------------
%	 
%----------------------------------------------------------------------------------------


\newpage 
\loadgeometry{marginlayout}  
\pagestyle{scrheadings} % Print headers again  
\KOMAoptions{
	headwidth=textwithmarginpar,
	headsepline=true,
	headsepline= 0.5pt,
	footwidth=textwithmarginpar,
} 
\onehalfspacing
\section{Propriétés de la racine carrée} 

\begin{proof} les deux nombres $a$ et $-a$ ont un carré égal à $a^2$.
\setlength{\abovedisplayskip}{3pt}
\setlength{\belowdisplayskip}{3pt} 
\[
(a)^2=(-a)^2
\] 
Si $a>0$ alors $\sqrt{a^2}=a$.\\
Si $a<0$, alors $-a>0$ et $\sqrt{a^2}=-a$. 
\end{proof}
\begin{proposition} Pour tout $a\in\mathbb{R}$ 
	\setlength{\abovedisplayskip}{3pt}
	\setlength{\belowdisplayskip}{3pt} 
	\[
	\sqrt{a^2}=\abs{a}
	\]
\end{proposition} 
\begin{example} \hfill 
	\begin{enumerate}[label=\alph*),leftmargin=*]
		\item $\big(3^2\big)^{\numprint{0.5}}  =\sqrt{3^2}=  3 $
		\item $\big((-3)^2\big)^{\numprint{0.5}}  =\sqrt{(-3)^2} = \sqrt{9}= 3 = -(-3) $
	\end{enumerate} 
\end{example}
\marginnote{
	\begin{tBox}[leftline=false]
		Les formules   généralisent celles déjà connues pour des produit de puissances à exposant entiers :
		\setlength{\abovedisplayskip}{3pt}
		\setlength{\belowdisplayskip}{3pt}
		\[
		\big(a\times b\big)^{\numprint{0.5}}  = a^{\numprint{0.5}} \times b^{\numprint{0.5}} 
		\]
		% 	\left( \frac{a}{b}\right)^{\numprint{0.5} = \frac{a^{\numprint{0.5}}{b^{\numprint{0.5}} 
				\end{tBox}
			}[0cm]
			\begin{theorem}
				Pour tout réels positifs non nuls $a$ et $b> 0$ :
				\setlength{\abovedisplayskip}{3pt}
				\setlength{\belowdisplayskip}{3pt} 
				\begin{align*}   
					\sqrt{a \times b} &= \sqrt{a} \times \sqrt{b}
					&
					\sqrt{\frac{a}{b}}  &= \frac{\sqrt{a}}{\sqrt{b}}  %&& ( b\neq 0) %\\
					%		(a\times b)^{\frac{1}{2}}& = a^\frac{1}{2} \times b^\frac{1}{2}
					%		&
					%		\left( \frac{a}{b}\right)^{\frac{1}{2}}& = \frac{a^\frac{1}{2}}{b^\frac{1}{2}}
				\end{align*}
			\end{theorem}
			\begin{proof} \emph{de la 1\iere égalité, au programme}\hfill \\
			$\sqrt{a}\sqrt{b}$ est un nombre positif, dont le carré vaut $ab$ car : 
				\begin{enumerate}[label=\textbullet]
					\item $\sqrt{a}\sqrt{b}\geqslant 0$ car produit de $\sqrt{a}\geqslant 0$ et $\sqrt{b}\geqslant 0$.
					\item  $(\sqrt{a}\sqrt{b})^2=(\sqrt{a})^2(\sqrt{b})^2 = ab$
				\end{enumerate} 
				Or $\sqrt{ab}$ désigne l'\textbf{unique} nombre positif dont le carré vaut $ab$.
				Donc on a $\sqrt{a}\sqrt{b} = \sqrt{ab}$. 
				%	\begin{description}
					%		\item[$\bullet$]$\frac{\sqrt{a}}{\sqrt{b}}\geqslant 0$ car \dotfill % quotient de $\sqrt{a}\geqslant 0$ et $\sqrt{b}> 0$.
					%		\item[$\bullet$] $\left(\frac{\sqrt{a}}{\sqrt{b}}\right)^2=$ \dotfill %\frac{(\sqrt{a})^2}{(\sqrt{b})^2} = \frac{a}{b}$.
					%	\end{description}
				%	Or $\sqrt{\frac{a}{b}}$ est l'unique nombre \dotfill 
				%	
				%	\null \dotfill   
			\end{proof}
			\begin{proof}[À vous] Démontrer la 2\ieme égalité.
			\end{proof}
			

%----------------------------------------------------------------------------------------
%	EXERCICES
%----------------------------------------------------------------------------------------

\newpage 
\loadgeometry{widelayout}  
\pagestyle{scrheadings} % Print headers again  
\KOMAoptions{
	headwidth=textwithmarginpar,
	headsepline=true,
	headsepline= 0.5pt,
	footwidth=textwithmarginpar,
} 
\onehalfspacing
\subsection{Exercices : racines carrées}
\begin{example}[Carrés parfaits à connaitre]  de collège.\hfill 
	\begin{multicols}{5}
		\begin{spacing}{1.8}
		\begin{enumerate}[label={$\arabic*^2=$}]
			\item 
			\item 
			\item 
			\item 
			\item 
			\item 
			\item 
			\item 
			\item 
			\item 
			\item 
			\item 
			\item 
			\item 
			\item 
		\end{enumerate} 
	\end{spacing}
	\end{multicols}
\end{example} 

\begin{exercise}[\cancel{\faCalculator}] Compléter lorsque c'est possible les expressions suivantes.
	\begin{multicols}{4}
		\begin{spacing}{2} 
			\begin{enumerate}[label={}, leftmargin=*]   
				\item $\sqrt{36}=\dotfill $
				\item $-\sqrt{64}=\dotfill $
				\item $\sqrt{-9}=\dotfill $
				\item $\sqrt{100}=\dotfill $
				\item $\left( \ldots\ldots\right)^2=5 $
				\item $\sqrt{3}^2=\dotfill $ 
				\item $\sqrt{5}\times \sqrt{5}=\dotfill $
				\item $-\sqrt{10}^2=\dotfill $
				\item $\left(-\sqrt{12}\right)^2=\dotfill$
				\item $\left(10\sqrt{2}\right)^2= \dotfill $
				\item $\left(2\sqrt{10}\right)^2= \dotfill $
				\item $\left(7\sqrt{2}\right)^2= \dotfill $ 
			\end{enumerate}
		\end{spacing}
	\end{multicols}
\end{exercise}
 
\begin{exercise}[\cancel{\faCalculator}] Corriger les expressions suivantes en rajoutant ou supprimant les $\sqrt{}$ (in)appropriés.
	\begin{multicols}{2}
		\begin{spacing}{2} 
			\begin{enumerate}[label={}, leftmargin=*]   
				\item  $\sqrt{36}=\sqrt{6}$ \dotfill  
				\item $\sqrt{49}^2=\sqrt{7}$ \dotfill
				\item $\sqrt{5}^2=25$ \dotfill
				\item $\sqrt{9^2}=9$ \dotfill
				\item $\sqrt{4^{10}}=\sqrt{4}^5$\dotfill
				\item $\sqrt{8^6}=8^3$\dotfill
				\item $3^2\sqrt{5}^2=\sqrt{45}$\dotfill 
				\item $ \sqrt{18}^2=\sqrt{3}^2\sqrt{2}^2$\dotfill 
			\end{enumerate}
		\end{spacing}
	\end{multicols}	\vspace{-0.5em}
\end{exercise}
\begin{exercise}[\cancel{\faCalculator}] Simplifier les racines de carrés suivantes :
	\begin{multicols}{2}
		\begin{spacing}{2} 
			\begin{enumerate}[label={}, leftmargin=*]   
				\item $\sqrt{5^2}=\abs{5}=$	\dotfill
				\item $\sqrt{(-5)^2}=\abs{-5}=$\dotfill  
				\item $\sqrt{(3-\sqrt{2})^2}=\abs{3-\sqrt{2}}=$\dotfill
				\item $\sqrt{\left(-9\right)^2}=$\dotfill  
				\item $\sqrt{(1-\pi)^2}=$\dotfill
				\item $\sqrt{\left(2-\tfrac{5}{3}\right)^2}=$\dotfill 
				\item $\sqrt{\left(1-\tfrac{5}{3}\right)^2}=$\dotfill 
				\item $\sqrt{(1-\pi)^2}=$\dotfill 
				\item  $\sqrt{(\numprint{3.14}-\pi)^2}=$\dotfill 
				\item  $-\left(-\sqrt{3^2}\right)^2=$\dotfill 
			\end{enumerate}
		\end{spacing}
	\end{multicols}	\vspace{-0.5em}
\end{exercise} 
\begin{exercise} Complétez pour trouver le meilleur encadement de $x$.\vspace{-0.5em}
	\begin{spacing}{2.0}
		\begin{enumerate}[leftmargin=*,label={}] 
			\item L'expression $\sqrt{-2x}$ est définie lorsque 
			\dotfill $\geqslant 0$, donc   pour $x\in$\dotfill
			\item L'expression  $\sqrt{x-2}$ est définie lorsque 
			\dotfill $\geqslant 0$, donc  pour  $x\in$ \dotfill 
			\item Si $\sqrt{x^2}=x$ alors $x$\dotfill; si $\sqrt{x^2}=-x$ alors $x$ \dotfill
			\item Si $\sqrt{(x-2)^2}=x-2$ alors\dotfill ; si $\sqrt{(x-2)^2}=2-x$ alors \dotfill
			\item Si $\sqrt{(3-2x)^2}=3-2x$ alors\dotfill 
			\dotfill  
		\end{enumerate} 
	\end{spacing}
	\vspace{-0.5em}
\end{exercise} 
%----------------------------------------------------------------------------------------
%	 
%----------------------------------------------------------------------------------------
\newpage 
 \begin{eBox} Pour $a, b>0$, on a $\sqrt{ab}=\sqrt{a}\sqrt{b}$ et $\sqrt{\frac{a}{b}}   = \frac{\sqrt{a}}{\sqrt{b}}$
 \end{eBox}
\begin{example}[Simplifier : Écrire une expression avec le terme sous la racine le plus petit possible]
	
	\noindent\parbox{0.55\linewidth}{
		$\begin{WithArrows}
			A & = \rule{0pt}{1.6em}\sqrt{12} \Arrow{identifier le plus grand carré facteur de $12$} \\
			&=  \rule{0pt}{1.6em}\sqrt{4\times 3} \Arrow{utiliser $\sqrt{ab}=\sqrt{a}\sqrt{b}$}\\
			& = \rule{0pt}{1.6em} \sqrt{4}\sqrt{3}  \\
			& = \rule{0pt}{1.6em} 2\sqrt{3}  
		\end{WithArrows}$
	}	\hfill 
	\parbox{0.20\linewidth}{	
		$\begin{WithArrows}
			B & = \rule{0pt}{1.6em}\sqrt{75}  \\
			&=  \rule{0pt}{1.6em} \\
			& =  \rule{0pt}{1.6em}  \\
			& =  \rule{0pt}{1.6em} 
		\end{WithArrows}$ 	
	}\parbox{0.20\linewidth}{	
		$\begin{WithArrows}
			C & = \rule{0pt}{1.6em}2\sqrt{18}  \\
			&=  \rule{0pt}{1.6em} \\
			& = \rule{0pt}{1.6em}  \\
			& =  \rule{0pt}{1.6em} 
		\end{WithArrows}$ 
	}	   
\end{example}

\begin{exercise}[\cancel{\faCalculator} À vous] Complétez par des entiers, afin de rendre les égalités vraies.
	
	\begin{multicols}{2}
		\begin{spacing}{2.5}
		\begin{enumerate}[label={}, leftmargin=*]
\item  $\sqrt{8}=\sqrt{\begin{tikzpicture}
\node [rectangle,draw,dashed, minimum size=0.7cm] (C) at (2,0.1) {\qquad};\end{tikzpicture}}\sqrt{\begin{tikzpicture}
\node [rectangle,draw,dashed, minimum size=0.7cm] (C) at (2,0.1) {\qquad};\end{tikzpicture}}= \begin{tikzpicture}
\node [rectangle,draw,dashed, minimum size=0.7cm] (C) at (2,0.1) {\qquad};\end{tikzpicture}\sqrt{\begin{tikzpicture}
\node [rectangle,draw,dashed, minimum size=0.7cm] (C) at (2,0.1) {\qquad};\end{tikzpicture}}$
\item  $\sqrt{32}=\sqrt{\begin{tikzpicture}
\node [rectangle,draw,dashed, minimum size=0.7cm] (C) at (2,0.1) {\qquad};\end{tikzpicture}}\sqrt{\begin{tikzpicture}
\node [rectangle,draw,dashed, minimum size=0.7cm] (C) at (2,0.1) {\qquad};\end{tikzpicture}}= \begin{tikzpicture}
\node [rectangle,draw,dashed, minimum size=0.7cm] (C) at (2,0.1) {\qquad};\end{tikzpicture}\sqrt{\begin{tikzpicture}
\node [rectangle,draw,dashed, minimum size=0.7cm] (C) at (2,0.1) {\qquad};\end{tikzpicture}}$
\item  $\sqrt{27}=\sqrt{\begin{tikzpicture}
\node [rectangle,draw,dashed, minimum size=0.7cm] (C) at (2,0.1) {\qquad};\end{tikzpicture}}\sqrt{\begin{tikzpicture}
\node [rectangle,draw,dashed, minimum size=0.7cm] (C) at (2,0.1) {\qquad};\end{tikzpicture}}= \begin{tikzpicture}
\node [rectangle,draw,dashed, minimum size=0.7cm] (C) at (2,0.1) {\qquad};\end{tikzpicture}\sqrt{\begin{tikzpicture}
\node [rectangle,draw,dashed, minimum size=0.7cm] (C) at (2,0.1) {\qquad};\end{tikzpicture}}$
\item  $\sqrt{50}=\sqrt{\begin{tikzpicture}
\node [rectangle,draw,dashed, minimum size=0.7cm] (C) at (2,0.1) {\qquad};\end{tikzpicture}}\sqrt{\begin{tikzpicture}
\node [rectangle,draw,dashed, minimum size=0.7cm] (C) at (2,0.1) {\qquad};\end{tikzpicture}}= \begin{tikzpicture}
\node [rectangle,draw,dashed, minimum size=0.7cm] (C) at (2,0.1) {\qquad};\end{tikzpicture}\sqrt{\begin{tikzpicture}
\node [rectangle,draw,dashed, minimum size=0.7cm] (C) at (2,0.1) {\qquad};\end{tikzpicture}}$
\item  $\sqrt{48}=\sqrt{\begin{tikzpicture}
\node [rectangle,draw,dashed, minimum size=0.7cm] (C) at (2,0.1) {\qquad};\end{tikzpicture}}\sqrt{\begin{tikzpicture}
\node [rectangle,draw,dashed, minimum size=0.7cm] (C) at (2,0.1) {\qquad};\end{tikzpicture}}= \begin{tikzpicture}
\node [rectangle,draw,dashed, minimum size=0.7cm] (C) at (2,0.1) {\qquad};\end{tikzpicture}\sqrt{\begin{tikzpicture}
\node [rectangle,draw,dashed, minimum size=0.7cm] (C) at (2,0.1) {\qquad};\end{tikzpicture}}$
\item  $\sqrt{200}=\sqrt{\begin{tikzpicture}
\node [rectangle,draw,dashed, minimum size=0.7cm] (C) at (2,0.1) {\qquad};\end{tikzpicture}}\sqrt{\begin{tikzpicture}
\node [rectangle,draw,dashed, minimum size=0.7cm] (C) at (2,0.1) {\qquad};\end{tikzpicture}}= \begin{tikzpicture}
\node [rectangle,draw,dashed, minimum size=0.7cm] (C) at (2,0.1) {\qquad};\end{tikzpicture}\sqrt{\begin{tikzpicture}
\node [rectangle,draw,dashed, minimum size=0.7cm] (C) at (2,0.1) {\qquad};\end{tikzpicture}}$ 
\item    $\sqrt{288}=\sqrt{144}\sqrt{\begin{tikzpicture}
	\node [rectangle,draw,dashed, minimum size=0.7cm] (C) at (2,0.1) {\qquad};\end{tikzpicture}}= \begin{tikzpicture}
\node [rectangle,draw,dashed, minimum size=0.7cm] (C) at (2,0.1) {\qquad};\end{tikzpicture}\sqrt{\begin{tikzpicture}
	\node [rectangle,draw,dashed, minimum size=0.7cm] (C) at (2,0.1) {\qquad};\end{tikzpicture}}$
\item  $\sqrt{98}=\sqrt{\begin{tikzpicture}
\node [rectangle,draw,dashed, minimum size=0.7cm] (C) at (2,0.1) {\qquad};\end{tikzpicture}}\sqrt{\begin{tikzpicture}
\node [rectangle,draw,dashed, minimum size=0.7cm] (C) at (2,0.1) {\qquad};\end{tikzpicture}}= \begin{tikzpicture}
\node [rectangle,draw,dashed, minimum size=0.7cm] (C) at (2,0.1) {\qquad};\end{tikzpicture}\sqrt{\begin{tikzpicture}
\node [rectangle,draw,dashed, minimum size=0.7cm] (C) at (2,0.1) {\qquad};\end{tikzpicture}}$
\item $\sqrt{900}=$\dotfill 
\item $\sqrt{8100}=$\dotfill 
\item $\sqrt{2500}=$\dotfill  
\item $\sqrt{160}=$\dotfill 
\item  $\sqrt{49\times 10^{4}}= \dotfill$ 
\item    $\sqrt{\begin{tikzpicture}
\node [rectangle,draw,dashed, minimum size=0.7cm] (C) at (2,0.1) {\qquad};\end{tikzpicture}}=\sqrt{36}\sqrt{2}= \begin{tikzpicture}
\node [rectangle,draw,dashed, minimum size=0.7cm] (C) at (2,0.1) {\qquad};\end{tikzpicture}\sqrt{\begin{tikzpicture}
\node [rectangle,draw,dashed, minimum size=0.7cm] (C) at (2,0.1) {\qquad};\end{tikzpicture}}$
\item    $\sqrt{\begin{tikzpicture}
\node [rectangle,draw,dashed, minimum size=0.7cm] (C) at (2,0.1) {\qquad};\end{tikzpicture}}=\sqrt{\begin{tikzpicture}
\node [rectangle,draw,dashed, minimum size=0.7cm] (C) at (2,0.1) {\qquad};\end{tikzpicture}}\sqrt{\begin{tikzpicture}
\node [rectangle,draw,dashed, minimum size=0.7cm] (C) at (2,0.1) {\qquad};\end{tikzpicture}}= 2\sqrt{6}$
\item    $\sqrt{\begin{tikzpicture}
\node [rectangle,draw,dashed, minimum size=0.7cm] (C) at (2,0.1) {\qquad};\end{tikzpicture}}=\sqrt{\begin{tikzpicture}
\node [rectangle,draw,dashed, minimum size=0.7cm] (C) at (2,0.1) {\qquad};\end{tikzpicture}}\sqrt{\begin{tikzpicture}
\node [rectangle,draw,dashed, minimum size=0.7cm] (C) at (2,0.1) {\qquad};\end{tikzpicture}}= 15\sqrt{2}$
\item    $\sqrt{\begin{tikzpicture}
\node [rectangle,draw,dashed, minimum size=0.7cm] (C) at (2,0.1) {\qquad};\end{tikzpicture}}=\sqrt{\begin{tikzpicture}
\node [rectangle,draw,dashed, minimum size=0.7cm] (C) at (2,0.1) {\qquad};\end{tikzpicture}}\sqrt{\begin{tikzpicture}
\node [rectangle,draw,dashed, minimum size=0.7cm] (C) at (2,0.1) {\qquad};\end{tikzpicture}}= 8\sqrt{5}$ 
\item  $\sqrt{\numprint{0.25}}=\sqrt{\frac{25}{100}}= \frac{\sqrt{\begin{tikzpicture}
		\node [rectangle,draw,dashed, minimum size=0.7cm] (C) at (2,0.1) {\qquad};\end{tikzpicture}}}{\sqrt{\begin{tikzpicture}
		\node [rectangle,draw,dashed, minimum size=0.7cm] (C) at (2,0.1) {\qquad};\end{tikzpicture}}} =\frac{\begin{tikzpicture}
	\node [rectangle,draw,dashed, minimum size=0.7cm] (C) at (2,0.1) {\qquad};\end{tikzpicture}}{\begin{tikzpicture}
	\node [rectangle,draw,dashed, minimum size=0.7cm] (C) at (2,0.1) {\qquad};\end{tikzpicture}} $
\item  $\sqrt{\numprint{0.36}}=\sqrt{\frac{36}{100}} = \frac{\sqrt{\begin{tikzpicture}
		\node [rectangle,draw,dashed, minimum size=0.7cm] (C) at (2,0.1) {\qquad};\end{tikzpicture}}}{\sqrt{\begin{tikzpicture}
		\node [rectangle,draw,dashed, minimum size=0.7cm] (C) at (2,0.1) {\qquad};\end{tikzpicture}}} =\frac{\begin{tikzpicture}
	\node [rectangle,draw,dashed, minimum size=0.7cm] (C) at (2,0.1) {\qquad};\end{tikzpicture}}{\begin{tikzpicture}
	\node [rectangle,draw,dashed, minimum size=0.7cm] (C) at (2,0.1) {\qquad};\end{tikzpicture}} =\dotfill$
\item  $\sqrt{\numprint{1.96}}=\sqrt{\frac{\begin{tikzpicture}
		\node [rectangle,draw,dashed, minimum size=0.7cm] (C) at (2,0.1) {\qquad};\end{tikzpicture}}{\begin{tikzpicture}
		\node [rectangle,draw,dashed, minimum size=0.7cm] (C) at (2,0.1) {\qquad};\end{tikzpicture}}}= \frac{\sqrt{\begin{tikzpicture}
		\node [rectangle,draw,dashed, minimum size=0.7cm] (C) at (2,0.1) {\qquad};\end{tikzpicture}}}{\sqrt{\begin{tikzpicture}
		\node [rectangle,draw,dashed, minimum size=0.7cm] (C) at (2,0.1) {\qquad};\end{tikzpicture}}} =\frac{\begin{tikzpicture}
	\node [rectangle,draw,dashed, minimum size=0.7cm] (C) at (2,0.1) {\qquad};\end{tikzpicture}}{\begin{tikzpicture}
	\node [rectangle,draw,dashed, minimum size=0.7cm] (C) at (2,0.1) {\qquad};\end{tikzpicture}} =\dotfill$ 
	\item  $\sqrt{\numprint{0.16}}= 
	\dotfill $
\item  $\sqrt{\numprint{1.21}}=  
\dotfill $
\item  $\sqrt{\numprint{0.64}}= 
\dotfill $
\item  $ \sqrt{\frac{\begin{tikzpicture}
			\node [rectangle,draw,dashed, minimum size=0.7cm] (C) at (2,0.1) {\qquad};\end{tikzpicture}}{\begin{tikzpicture}
			\node [rectangle,draw,dashed, minimum size=0.7cm] (C) at (2,0.1) {\qquad};\end{tikzpicture}}}= \frac{\sqrt{\begin{tikzpicture}
			\node [rectangle,draw,dashed, minimum size=0.7cm] (C) at (2,0.1) {\qquad};\end{tikzpicture}}}{\sqrt{\begin{tikzpicture}
			\node [rectangle,draw,dashed, minimum size=0.7cm] (C) at (2,0.1) {\qquad};\end{tikzpicture}}} = \frac{2}{3}  $	
\item  $ \sqrt{\frac{\begin{tikzpicture}
			\node [rectangle,draw,dashed, minimum size=0.7cm] (C) at (2,0.1) {\qquad};\end{tikzpicture}}{\begin{tikzpicture}
			\node [rectangle,draw,dashed, minimum size=0.7cm] (C) at (2,0.1) {\qquad};\end{tikzpicture}}}= \frac{\sqrt{\begin{tikzpicture}
			\node [rectangle,draw,dashed, minimum size=0.7cm] (C) at (2,0.1) {\qquad};\end{tikzpicture}}}{\sqrt{\begin{tikzpicture}
			\node [rectangle,draw,dashed, minimum size=0.7cm] (C) at (2,0.1) {\qquad};\end{tikzpicture}}} = \frac{5}{7}  $
\item  $ \sqrt{\frac{\begin{tikzpicture}
		\node [rectangle,draw,dashed, minimum size=0.7cm] (C) at (2,0.1) {\qquad};\end{tikzpicture}}{\begin{tikzpicture}
		\node [rectangle,draw,dashed, minimum size=0.7cm] (C) at (2,0.1) {\qquad};\end{tikzpicture}}}= \frac{\sqrt{\begin{tikzpicture}
		\node [rectangle,draw,dashed, minimum size=0.7cm] (C) at (2,0.1) {\qquad};\end{tikzpicture}}}{\sqrt{\begin{tikzpicture}
		\node [rectangle,draw,dashed, minimum size=0.7cm] (C) at (2,0.1) {\qquad};\end{tikzpicture}}} = \frac{5^2}{7^3}  $ 
		\end{enumerate}
	\end{spacing}
	\end{multicols} 
\end{exercise}
 
\begin{exercise}[\cancel{\faCalculator}]\label{chap08:exo:radicaux:simplifier} Simplifier   en rendant  le terme sous la racine le plus petit possible.
	\begin{multicols}{2} 
		\begin{spacing}{2.0}
			\begin{enumerate}[label={}, leftmargin=*] 
		\item $3\sqrt{50}=3\times \sqrt{25}\sqrt{2}=15\sqrt{2}$
		\item $2\sqrt{48}=\dotfill$ 
		\item $5\sqrt{63}=\dotfill$ 
		\item $10\sqrt{8}=\dotfill$
		\item $2\sqrt{12}=\dotfill$
		\item $3\sqrt{20}=\dotfill$ 
		\item $4\sqrt{27}=\dotfill$ 
		\item $5\sqrt{18}=\dotfill$ 
		\item $5\sqrt{200}=\dotfill$
		\item $3\sqrt{125}=\dotfill$
		\item $3\sqrt{80}=\dotfill$
		\item $10\sqrt{96}=\dotfill$
		\item $7\sqrt{88}=\dotfill$ 
		\item $10\sqrt{75}=\dotfill$ 
		%			\item $\sqrt{32}$ 
	\end{enumerate}
	\end{spacing}
	\end{multicols}
\end{exercise}  
\begin{example}[Simplifier des produits ou quotients de radicaux] \hfill 
 	
 	\noindent\parbox{0.60\linewidth}{
 		$\begin{WithArrows}
 			A & = \sqrt{2}\times 3 \Arrow{mettre la racine après le coefficient}\\
 			& = 3\sqrt{2}\\
 			B & = 3\sqrt{2}\times \sqrt{3} \Arrow{utiliser l'identité $\sqrt{a}\sqrt{b}=\sqrt{ab}$}\\
 			& = 3\sqrt{6}
 		\end{WithArrows}$
 	}\hfill
 	\parbox{0.35\linewidth}{
 		$\begin{WithArrows}
 			C & = 2\sqrt{3} \times  5\sqrt{3} \\
 			&=    \\
 			& =   
 		\end{WithArrows}$
 	}
 \end{example} 
 \begin{exercise}[\cancel{\faCalculator}]\label{chap08:exo:radicaux:simplifier2}\hfill\vspace{-0.5em} 
 	\begin{multicols}{2}
 		\begin{spacing}{2.0}
 		\begin{enumerate}[label={}, leftmargin=*]  
 			\item $\sqrt{7}\times \sqrt{5}=\dotfill$
 			\item $\sqrt{5}\times 5=\dotfill$
 			\item $\sqrt{5}\times \sqrt{2}=\dotfill$
 			\item $\sqrt{11}\times2 \times  \sqrt{3}=\dotfill$
 			\item $5\sqrt{2} \times  \sqrt{2}=\dotfill$
 			\item $\sqrt{5}\times \sqrt{45}=\dotfill$
 			\item $5\sqrt{21}\times 2\sqrt{3}=\dotfill$  
 			\item $2\sqrt{3} \times  \sqrt{27}=\dotfill$
 			\item $(\sqrt{7})^3=\dotfill$
 			\item $\sqrt{3^8}=\sqrt{\big(3^4\big)^2}=\dotfill$
 			\item $\sqrt{4^5}=\sqrt{4^4\times 4}=\dotfill$
 			\item $\sqrt{5^{10}}=\dotfill$
 			\item $\sqrt{36^3}= \dotfill$
 			\item $(\sqrt{6})^4=\dotfill$
 			\item $(6\sqrt{5})^2=\dotfill$
 			\item $(7\sqrt{2})^2=\dotfill$ 
 			\item $(2\sqrt{5})^3=\dotfill$ 
 			\item $\sqrt{10^{-9}}\times \sqrt{10^{15}}=\dotfill$
 			\item $\sqrt{5^{7}}\times \sqrt{5^{-4}}=\dotfill$
 			\item $\sqrt{5^{-7}}\times \sqrt{5^{-4}}=\dotfill$
 		\end{enumerate}
 	\end{spacing}
 	\end{multicols}
 \end{exercise}
\begin{exercise}[Comparer]  
	Écrire sous forme $\sqrt{k}$ ou $k\in\mathbb{N}$. Quelle est la médiane de la série de nombres ?\vspace{-0.5em}  
	\begin{multicols}{4}
		\begin{spacing}{2.0}
		\begin{enumerate}[label={}, leftmargin=*] 
			\item $8\sqrt{2}=\dotfill $
			\item $7\sqrt{3}=\dotfill $
			\item $5\sqrt{5}=\dotfill $
			\item $5\sqrt{6}=\dotfill $ 
			\item $4\sqrt{10}=\dotfill $
			\item $3\sqrt{17}=\dotfill $ 
			\item $2\sqrt{37}=\dotfill $
		\end{enumerate}
	\end{spacing}
	\end{multicols}
	
\end{exercise} 

%----------------------------------------------------------------------------------------
%	 
%----------------------------------------------------------------------------------------
\newpage 
\begin{example}[Réduire des sommes de radicaux] Réduire les expressions suivantes :
	
	\noindent\parbox{0.55\linewidth}{
		$\begin{WithArrows}
			A & = \rule{0pt}{1.3em}\sqrt{2} + 3\sqrt{2}\Arrow{comme pour réduire $x+3x$} \\
			&=\rule{0pt}{1.3em} 4\sqrt{2}   \\
			B & =\rule{0pt}{1.3em} \sqrt{27}+5\sqrt{3} \Arrow{simplifier les racines} \\
			& = \rule{0pt}{1.3em}\sqrt{9}\sqrt{3}+5\sqrt{3} \\
			& =  3\sqrt{3}+5\sqrt{3}  =   8\sqrt{3}\\
			C & = \rule{0pt}{1.3em} 4\sqrt{2}-7\sqrt{5}+~ \sqrt{5} \\
			&=  \rule{0pt}{1.4em}    
		\end{WithArrows}$ 
	}
	\hfill 
	\parbox{0.40\linewidth}{
		$\begin{WithArrows}
			D & = \rule{0pt}{1.3em}5\sqrt{2}+3-2\sqrt{2}+5 \\
			& =  \rule{0pt}{1.3em} \\
			E & = \rule{0pt}{1.3em}2\sqrt{27}-5\sqrt{12} \\
			&= \rule{0pt}{1.3em} \\
			&=  \rule{0pt}{1.3em}\\
			&= \rule{0pt}{1.3em} 
		\end{WithArrows}$ 
	}
\end{example}   
\begin{exercise}[\cancel{\faCalculator}] \label{chap08:exo:radicaux:reduire} Mêmes consignes :\hfill \vspace{-0.5em}  
	\begin{multicols}{2}\begin{spacing}{2}
		\begin{enumerate}[label={},leftmargin=*]  
			\item $\sqrt{3}+\sqrt{3}=\dotfill$
			\item $6\sqrt{7}+3\sqrt{7}=\dotfill$
			\item $10\sqrt{3}-\sqrt{3}=\dotfill$
			\item $\sqrt{2}+5\sqrt{3}-2\sqrt{2}+6\sqrt{3}=\dotfill$
			\item $\sqrt{27}-2\sqrt{3}=\dotfill$ 
			\item $9\sqrt{7}+\sqrt{63}=\dotfill$  
			\item $7\sqrt{8}+3\sqrt{2}=\dotfill$
			\item $2\sqrt{80}-3\sqrt{20}=\dotfill$
			\item $2\sqrt{2}+8-\sqrt{2}+1=\dotfill$
			\item $7\sqrt{3}+\sqrt{3}-2=\dotfill$
			\item $10+4\sqrt{7}-2+\sqrt{7}=\dotfill$
			\item $\sqrt{2}+\sqrt{3}-2\sqrt{2}=\dotfill$ 
		\end{enumerate}\end{spacing}
	\end{multicols}
\end{exercise}

%\newpage
\begin{example}[Je fais] Développer simplifier et réduire les produits de radicaux :
	
	\noindent \parbox{1\linewidth}{	
		$\begin{WithArrows}
			A & =\rule{0pt}{1.3em}(\sqrt{3}-2)(\sqrt{6}+5)  = \rule{0pt}{1.3em} \\
			B & =\rule{0pt}{1.3em}(\sqrt{8}+3)(\sqrt{2}-1)    = \rule{0pt}{1.3em}  
		\end{WithArrows}$ 
	}
\end{example}
\begin{exercise}[\cancel{\faCalculator}] \label{chap08:exo:radicaux:developper} Même consignes \vspace{-0.5em}  
%	\begin{multicols}{2}
		\begin{spacing}{2}
			\begin{enumerate}[label={},leftmargin=*]  
				\item $3(4+\sqrt{2})=\dotfill$  $\sqrt{3}(4+\sqrt{3})=\dotfill$
				\item $\sqrt{3}(\sqrt{12}+2\sqrt{3})=\dotfill$
				\item$(\sqrt{7}+3)(\sqrt{3}+5)=\dotfill$
				\item$(3+\sqrt{11})(3-\sqrt{11})=\dotfill$ 
				\item$(3-2\sqrt{5})^2=\dotfill$   
				\item$(\sqrt{3}-\sqrt{2})(\sqrt{3}+\sqrt{2})=\dotfill$
				\item$(\sqrt{5}-\sqrt{3})(\sqrt{3}+2)=\dotfill$
				\item$(\sqrt{8}-\sqrt{12})(\sqrt{2}-\sqrt{3})=\dotfill$
				\item$(4\sqrt{5}-\sqrt{2})(3\sqrt{2}-\sqrt{5})=\dotfill$
		\end{enumerate}
	\end{spacing}
%	\end{multicols} 
\end{exercise} 

\begin{example} Transformer les fractions suivantes en une fraction égale ayant un dénominateur entier : \hfill 
	
	\noindent\hfill\parbox{0.30\linewidth}{	
		$\begin{WithArrows}
			\frac{7}{3\sqrt{5}}& = \frac{7\times \sqrt{5} }{3\sqrt{5}\times \sqrt{5}}  \\
			&= \frac{7\sqrt{5}}{\sqrt{25}} \\
			& = \frac{7\sqrt{5}}{5}  \\ 
			& \rule{0pt}{1.2em} %= \rule{0pt}{1.1em}  
		\end{WithArrows}$ 	
	} \hfill\parbox{0.30\linewidth}{	
		$\begin{WithArrows}
			\frac{3}{\sqrt{12}} &=   \frac{3\times\sqrt{12}}{\sqrt{12}\times\sqrt{12}} \\
			&=  \frac{3\sqrt{12}}{12} \\
			& =  \frac{3\sqrt{4}\sqrt{3}}{12}  \\
			& =   \frac{6\sqrt{3}}{12}= \frac{\sqrt{3}}{2} 
		\end{WithArrows}$ 	
	}  \hfill\parbox{0.30\linewidth}{	
		$\begin{WithArrows}
			\frac{4}{\sqrt{8}} &=  \frac{4\times\sqrt{8}}{\sqrt{8}\times\sqrt{8}} \\
			&=  \frac{4\times\sqrt{4}\sqrt{2}}{8}\\
			& = \frac{8\sqrt{2}}{8}  \\
			& =  \sqrt{2}
		\end{WithArrows}$
	}  
\end{example} 
\begin{exercise}[\cancel{\faCalculator} À vous]\label{chap08:exo:identite:radicaux1} Mêmes consignes. \vspace{-0.5em} 
	\begin{multicols}{2}
		\begin{spacing}{2}
		\begin{enumerate}[label={},leftmargin=*]
			\item $\frac{5}{\sqrt{7}}=\dotfill$
			\item $\frac{10}{\sqrt{6}}=\dotfill$
			\item $\frac{2}{\sqrt{10}}=\dotfill$
			\item $\frac{\sqrt{3}}{\sqrt{5}}=\dotfill$
			\item $\frac{\sqrt{3}}{\sqrt{7}}=\dotfill$
			\item $\frac{-\sqrt{9}}{\sqrt{3}}=\dotfill$
			\item $\frac{14}{3\sqrt{7}}=\dotfill$
			\item $\frac{7}{2\sqrt{3}}=\dotfill$
			\item $\frac{\sqrt{5}}{\sqrt{10}}=\dotfill$
			\item $\frac{\sqrt{18}}{2\sqrt{2}}=\dotfill$
		\end{enumerate}
	\end{spacing}
	\end{multicols}
\end{exercise}

\begin{example}  Mêmes consignes. \hfill 
	
	\noindent\hfill\parbox{0.32\linewidth}{	
		$\begin{WithArrows}
			\frac{3}{2+\sqrt{3}}& = \frac{3(2-\sqrt{3})}{(2+\sqrt{3})(2-\sqrt{3})} \\
			&= \frac{6-3\sqrt{3})}{2^2-\sqrt{3}^2}  \\
			& =\frac{6-3\sqrt{3})}{4-3} \\
			& = 6-3\sqrt{3}
		\end{WithArrows}$ 	
	} \hfill\parbox{0.32\linewidth}{	
		$\begin{WithArrows}
			\frac{4}{\sqrt{5}-2} &=    \frac{4(\sqrt{5}+2)}{(\sqrt{5}-2)(\sqrt{5}+2)} \\
			&=  \frac{4\sqrt{5}+8}{\sqrt{5}^2-2^2} \\
			& =   \frac{4\sqrt{5}+8}{5-4}  \\
			& =   4\sqrt{5}+8
		\end{WithArrows}$ 	
	}  \hfill\parbox{0.32\linewidth}{	
		$\begin{WithArrows}
			\frac{6}{\sqrt{5}-\sqrt{2}} &=  \frac{6(\sqrt{5}+\sqrt{2})}{(\sqrt{5}-\sqrt{2})(\sqrt{5}+\sqrt{2})} \\
			&=  \frac{6\sqrt{5}+6\sqrt{2}}{\sqrt{5}^2-\sqrt{2}^2}\\
			& = \frac{6\sqrt{5}+6\sqrt{2}}{3}  \\
			& =   2\sqrt{5}+2\sqrt{2}
		\end{WithArrows}$
	}  
\end{example} 
\begin{exercise}[\cancel{\faCalculator}] \label{chap08:exo:identite:radicaux2} Mêmes consignes.
	\vspace{-0.5em} 
%	\begin{multicols}{1} 
		\begin{spacing}{2}
		\begin{enumerate}[label={},leftmargin=*]
			\item  $\frac{2}{4+\sqrt{3}}=\dotfill$
			\item $\frac{5}{7-\sqrt{5}}=\dotfill$
			\item $\frac{10}{5+\sqrt{2}}=\dotfill$
			\item $\frac{15}{\sqrt{7}-5}=\dotfill$
			\item $\frac{4}{3\sqrt{2}-10}=\dotfill$
			\item $\frac{3}{7-4\sqrt{3}}=\dotfill$
			\item $\frac{4\sqrt{2}}{\sqrt{2}-1}=\dotfill$
			\item $\frac{\sqrt{3}-\sqrt{2}}{\sqrt{3}+\sqrt{2}}=\dotfill$
		\end{enumerate}
	\end{spacing}
%	\end{multicols}
\end{exercise}


%----------------------------------------------------------------------------------------
%	 
%---------------------------------------------------------------------------------------- 
\newpage  
\begin{tBox}
\textbf{Bilan à retenir} Pour rendre rationnel le dénominateur d'une fraction :
	\begin{enumerate}[label=\textbullet, leftmargin=*]
		\item si le dénominateur est un produit ayant pour facteur $\sqrt{a}$, alors on multiplie par   $\sqrt{a}$.
		\item si le dénominateur est de la forme $a\pm\sqrt{b}$  ( ou $\sqrt{a}\pm\sqrt{b}$ ) alors on multipliera par le conjugué.
	\end{enumerate}
\end{tBox}
\begin{exercise} Complétez.\vspace{-0.5em}
	\begin{spacing}{2.0}
		\begin{enumerate}[leftmargin=*,label={}]  
			\item Si $x\ldots 0$ et $y\ldots 0$  alors $\sqrt{xy}=\sqrt{x}\sqrt{y}$ et  $\sqrt{\frac{x}{y}}=\frac{\sqrt{x}}{\sqrt{y}}$ 
			\item Si $a<0$ alors $\sqrt{16a^2}=$\dotfill
			\item Si $a<0$ et $b\geqslant 0$ alors $\sqrt{a^2b}=$\dotfill  
			\item Si $a>0$ et $b\geqslant 0$ alors  $\sqrt{\frac{4b^2}{9a^2}}=$\dotfill  
			\item Si $a>0$  et $b<0$ alors  $\sqrt{\tfrac{25b^4}{36a^2}}=$\dotfill  
%			\item Si $a>0$ et $b\geqslant 0$ alors  $\sqrt{\frac{4b^2}{9a^2}}=$\dotfill  
%			\item L'égalité $\sqrt{x^2-1}=\sqrt{x-1}\sqrt{x+1}$ est vraie lorsque \dotfill   
		\end{enumerate} 
	\end{spacing}
	\vspace{-0.5em}
\end{exercise}    
\begin{problem}  [classique] \hfill 
\begin{enumerate}[label={\arabic*)},leftmargin=*]
	\item Développer simplifier et réduire $(3+5\sqrt{2})^2$.  En déduire une écriture simplifiée de $\sqrt{59+30\sqrt{2}}$.
	\item Développer simplifier et réduire  $(2-\sqrt{5})^2$. En déduire une écriture simplifiée de $\sqrt{9-4\sqrt{5}}$.
\end{enumerate} 
\end{problem}   
\begin{problem}\hfill  \\
	\noindent
	\parbox{0.7\linewidth}{
	\begin{enumerate}[leftmargin=*,label=\arabic*)]
%		\item Compléter par $<$, $>$ ou $=$ : 
%		\begin{multicols}{4}
%			\begin{enumerate}[label={\textbullet}, leftmargin=*]   
%				\item[] $\sqrt{10^2+3^2}\ldots 10 + 3 $  
%				\item[] $\sqrt{121+25}\ldots 11 + 5  $  
%				\item[] $\sqrt{9-4}\ldots \sqrt{9}-\sqrt{4}$
%				\item[] $\sqrt{8^2-6^2}\ldots 8 -6$  
%			\end{enumerate}
%		\end{multicols}	
		\item À l'aide du théorème de Pythagore, exprimer la longueur de l'hypoténuse à l'aide de $a$ et $b$.
		\item Expliquer pourquoi cette figure illustre l'inégalité pour tout $a$ et $b>0$, on a $\sqrt{a+b}<\sqrt{a}+\sqrt{b}$.
	\end{enumerate} 
}
\hfill 
\parbox{0.25\linewidth}{
%		\begin{figure*}[hb!] 
	\begin{tikzpicture}[scale=.8]
		%			\tkzInit[xmin=-1.0,ymin=-1.0,xmax=6.5,ymax=2.8]
		%			\tkzClip
		\tkzDefPoint(0,0){A} \tkzDefPoint(5,0){C}
		\tkzDefTriangle[two angles = 27 and 90](A,C)\tkzGetPoint{B}
		\tkzDrawPolygon[line width=1.2pt](A,B,C)
		% 			\tkzLabelSegment[below](A,C){$2\sqrt{xy}$}
		\tkzLabelSegment[below](A,C){$\sqrt{a}$}
		\tkzLabelSegment[right](B,C){$\sqrt{b}$}
		\tkzLabelSegment[above left](A,B){$?$}
		%    	\tkzDrawPoints[color=red](A,B,C)
		%   	 \tkzLabelPoints[below](A,C)
		%    		\tkzLabelPoints[above right](B)
		\tkzMarkRightAngle[  size=0.4](B,C,A)
	\end{tikzpicture} 
	%	\caption{Si $x$ et $y\geqslant 0$ alors  $x+y\geqslant 2\sqrt{xy}$} 
	%		\end{figure*}  
}
\end{problem}

%\begin{problem} Si dessous une suite de figures :
%	\begin{figure*}[hb!] 
%		\begin{tikzpicture}[scale=.8]
%			%			\tkzInit[xmin=-1.0,ymin=-1.0,xmax=6.5,ymax=2.8]
%			%			\tkzClip
%			\tkzDefPoint(0,0){A} \tkzDefPoint(2,0){C}
%			\tkzDefTriangle[two angles = 45 and 90](A,C)\tkzGetPoint{B}
%			\tkzDrawPolygon[line width=1.2pt](A,B,C)
%			% 			\tkzLabelSegment[below](A,C){$2\sqrt{xy}$}
%			\tkzLabelSegment[below](A,C){$3\sqrt{6}$}
%			\tkzLabelSegment[right](B,C){$3\sqrt{6}$}
%			\tkzLabelSegment[above left](A,B){$a$}
%			%    	\tkzDrawPoints[color=red](A,B,C)
%			%   	 \tkzLabelPoints[below](A,C)
%			%    		\tkzLabelPoints[above right](B)
%			\tkzMarkRightAngle[  size=0.4](B,C,A)
%			\begin{scope}[xshift=4cm]
%				\tkzDefPoint(0,0){A} \tkzDefPoint(3,0){C}
%				\tkzDefTriangle[two angles = 45 and 90](A,C)\tkzGetPoint{B}
%				\tkzDrawPolygon[line width=1.2pt](A,B,C)
%				% 			\tkzLabelSegment[below](A,C){$2\sqrt{xy}$}
%				\tkzLabelSegment[below](A,C){$5\sqrt{10}$}
%				\tkzLabelSegment[right](B,C){$5\sqrt{10}$}
%				\tkzLabelSegment[above left](A,B){$b$}
%				%    	\tkzDrawPoints[color=red](A,B,C)
%				%   	 \tkzLabelPoints[below](A,C)
%				%    		\tkzLabelPoints[above right](B)
%				\tkzMarkRightAngle[  size=0.4](B,C,A)
%			\end{scope}
%			\begin{scope}[xshift=9cm]
%				\tkzDefPoint(0,0){A} \tkzDefPoint(4,0){C}
%				\tkzDefTriangle[two angles = 45 and 90](A,C)\tkzGetPoint{B}
%				\tkzDrawPolygon[line width=1.2pt](A,B,C)
%				% 			\tkzLabelSegment[below](A,C){$2\sqrt{xy}$}
%				\tkzLabelSegment[below](A,C){$7\sqrt{14}$}
%				\tkzLabelSegment[right](B,C){$7\sqrt{14}$}
%				\tkzLabelSegment[above left](A,B){$c$}
%				%    	\tkzDrawPoints[color=red](A,B,C)
%				%   	 \tkzLabelPoints[below](A,C)
%				%    		\tkzLabelPoints[above right](B)
%				\tkzMarkRightAngle[  size=0.4](B,C,A)
%			\end{scope}
%			\begin{scope}[xshift=15.5cm]
%				\tkzDefPoint(0,0){A} \tkzDefPoint(5,0){C}
%				\tkzDefTriangle[two angles = 45 and 90](A,C)\tkzGetPoint{B}
%				\tkzDrawPolygon[line width=1.2pt](A,B,C)
%				% 			\tkzLabelSegment[below](A,C){$2\sqrt{xy}$}
%				\tkzLabelSegment[below](A,C){$\ldots\sqrt{\ldots}$}
%				\tkzLabelSegment[right](B,C){$\ldots\sqrt{\ldots}$}
%				\tkzLabelSegment[above left](A,B){$\ldots $}
%				%    	\tkzDrawPoints[color=red](A,B,C)
%				%   	 \tkzLabelPoints[below](A,C)
%				%    		\tkzLabelPoints[above right](B)
%				\tkzMarkRightAngle[  size=0.4](B,C,A)
%			\end{scope}
%		\end{tikzpicture} 
%		%	\caption{Si $x$ et $y\geqslant 0$ alors  $x+y\geqslant 2\sqrt{xy}$} 
%	\end{figure*}  
%	\begin{enumerate}[label=\arabic*),leftmargin=*]
%		\item 	Trouver les longueurs manquantes pour les 3 premières, et exprimer les à l'aide de radicaux.
%		\item Quelle conjecture pouvez vous faire pour la suite des figures ? Démontrer la.
%	\end{enumerate} 
%\end{problem}
 
%----------------------------------------------------------------------------------------
%	 
%----------------------------------------------------------------------------------------
\begin{proof}[solution de l'exercice \ref{chap08:exo:radicaux:simplifier}] 
	$a=8\sqrt{3}$; $b=15\sqrt{7}$; $c=20\sqrt{2}$; $d=6\sqrt{5}$; $e=12\sqrt{3}$; $f=50\sqrt{2}$; $g=15\sqrt{5}$; $h=12\sqrt{5}$ ;  $i=40\sqrt{6}$ ;  $j=14\sqrt{22}$. 
\end{proof}
\begin{proof}[solution de l'exercice \ref{chap08:exo:radicaux:simplifier2}] 
 $a=\sqrt{35}$; $b=5\sqrt{5}$; $c=\sqrt{10}$; $d=2\sqrt{33}$; $e=10$; $f=15$; $g=20\sqrt{7}$;
 $h=18$;
 $h=7\sqrt{7}$;
 $h=3^4$;
 $i=32$;
 $j=5^5$; 
 $k=6^3$; $l=36$; $m=36\times 5$;$n=98$; $p=40$; $q=1000$; $r=5\sqrt{5}$; $s=5^5\sqrt{5}$
\end{proof} 
\begin{proof}[solution de l'exercice \ref{chap08:exo:radicaux:reduire}]
\begin{pycode}
from re import sub
import sympy as sp

x, y, z, t = sp.symbols('x y z t')
k, m, n = sp.symbols('k m n', integer=True)
f, g, h = sp.symbols('f g h', cls=sp.Function)

# Create a list of functions to include in the table
funcs = [  '2*sp.sqrt(2)',
'9*sp.sqrt(7)',
'9*sp.sqrt(3)',
'11*sp.sqrt(3)-1*sp.sqrt(2)',
'sp.sqrt(27)-2*sp.sqrt(3)',
'9*sp.sqrt(7)+1*sp.sqrt(63)',
'7*sp.sqrt(8)+3*sp.sqrt(2)', 
'2*sp.sqrt(80)-3*sp.sqrt(20)', 
'2*sp.sqrt(2)-1*sp.sqrt(2)+1', 
'7*sp.sqrt(3)+1*sp.sqrt(3)-2', 
'8+5*sp.sqrt(7)', 
'1*sp.sqrt(3)-1*sp.sqrt(2)',
]
n=0
L=['A', 'B', 'C', 'D', 'E', 'F', 'G', 'H', 'I', 'J','K', 'L', 'M',  ]
for func in funcs:
	# Put in some vertical space when switching to arc and hyperbolic funcs
	developpee =  eval('sp.expand('+ func +')')
	print( "$"+L[n] + "  =" + sp.latex( eval('developpee') )  + r'$; ' )
	n     = n + 1
\end{pycode}
\end{proof}
% 
\begin{proof}[solution de l'exercice \ref{chap08:exo:radicaux:developper}] 
	\begin{pycode}
from re import sub
import sympy as sp

x, y, z, t = sp.symbols('x y z t')
k, m, n = sp.symbols('k m n', integer=True)
f, g, h = sp.symbols('f g h', cls=sp.Function)

# Create a list of functions to include in the table
funcs = [  '3* (4+sp.sqrt(2))'
,'sp.sqrt(3)*(4+sp.sqrt(3))'
,'sp.sqrt(3)*(sp.sqrt(12)+2*sp.sqrt(3))'
,'(sp.sqrt(7)+3)*(sp.sqrt(3)+5)'
,'(3+sp.sqrt(11))*(3-sp.sqrt(11))'  
,'(3-2*sp.sqrt(11))**2'  
,'(sp.sqrt(3)-sp.sqrt(2))*(sp.sqrt(3)+sp.sqrt(2))'
,'(sp.sqrt(5)-sp.sqrt(3))*(sp.sqrt(3)+2)'
,'(sp.sqrt(8)-sp.sqrt(12))*(sp.sqrt(2)-sp.sqrt(3))'
,'(4*sp.sqrt(5)-sp.sqrt(2))*(3*sp.sqrt(2)-sp.sqrt(5))'
]
n=0
L=['A', 'B', 'C', 'D', 'E', 'F', 'G', 'H', 'I', 'J'  ]
for func in funcs:
	# Put in some vertical space when switching to arc and hyperbolic funcs
	developpee =  eval('sp.expand('+ func +')')
	print( "$"+L[n] + "  =" + sp.latex( eval('developpee') )  + r'$; ' )
	n     = n + 1
\end{pycode}
\end{proof} 




\begin{proof}[solution de l'exercice \ref{chap08:exo:identite:radicaux1}]\scriptsize
\begin{pycode}
from re import sub
import sympy as sp

x, y, z, t = sp.symbols('x y z t')
k, m, n = sp.symbols('k m n', integer=True)
f, g, h = sp.symbols('f g h', cls=sp.Function)

# Create a list of functions to include in the table
funcs = ['5/(sp.sqrt(7))',
'10/(sp.sqrt(6))',
'2/(sp.sqrt(10))',
'(sp.sqrt(3))/(sp.sqrt(5))',
'(sp.sqrt(3))/(sp.sqrt(7))',
'(-sp.sqrt(9))/(sp.sqrt(3))',
'(14)/(3*sp.sqrt(7))',
'(7)/(2*sp.sqrt(3))',
'(sp.sqrt(5))/(sp.sqrt(10))',
'(sp.sqrt(18))/(2*sp.sqrt(2))'
]
n=0
L=['A', 'B', 'C', 'D', 'E', 'F', 'G', 'H', 'I' , 'J' ]
for func in funcs:
	# Put in some vertical space when switching to arc and hyperbolic funcs
	simplifiee =  eval('sp.radsimp('+ func +')')
	print( "$"+L[n] + " =" + sp.latex( eval('simplifiee') )  + r'$; ' )
	n     = n + 1
\end{pycode}
\end{proof}	
\begin{proof}[solution de l'exercice \ref{chap08:exo:identite:radicaux2}]\scriptsize
	\begin{pycode}
from re import sub
import sympy as sp

x, y, z, t = sp.symbols('x y z t')
k, m, n = sp.symbols('k m n', integer=True)
f, g, h = sp.symbols('f g h', cls=sp.Function)

# Create a list of functions to include in the table
funcs = ['2/(4+sp.sqrt(3))',
'5/(7-sp.sqrt(5))',
'10/(5+sp.sqrt(2))',
'15/( sp.sqrt(7)-5)',
'4/(3*sp.sqrt(2)-10)',
'3/(7-4*sp.sqrt(3))',
'(4*sp.sqrt(2))/(sp.sqrt(2)-1)',
'(sp.sqrt(3)-sp.sqrt(2))/(sp.sqrt(3)+sp.sqrt(2))' 
]
n=0
L=['A', 'B', 'C', 'D', 'E', 'F', 'G', 'H', 'I'  ]
for func in funcs:
	# Put in some vertical space when switching to arc and hyperbolic funcs
	simplifiee =  eval('sp.radsimp('+ func +')')
	print( "$"+L[n] + " =" + sp.latex( eval('simplifiee') )  + r'$; ' )
	n     = n + 1
\end{pycode}
\end{proof}


%----------------------------------------------------------------------------------------
%	EXERCICES
%----------------------------------------------------------------------------------------

\newpage 
\loadgeometry{widelayout}  
\pagestyle{scrheadings} % Print headers again  
\KOMAoptions{
	headwidth=textwithmarginpar,
	headsepline=true,
	headsepline= 0.5pt,
	footwidth=textwithmarginpar,
} 
\onehalfspacing
\section{Club maths : puissances et multiplications de radicaux}

\begin{exercise} Pour chaque cas trouver des nombres entiers $a$, $b$ ou $c$ qui rendent l'égalité vraie.\vspace{-0.5em}
\begin{multicols}{2} 
	\begin{spacing}{2}
	\begin{enumerate}[label={\arabic*)},leftmargin=*]
		\item  $(a+\sqrt{b})(a-\sqrt{b})=3$
		\item  $(a+\sqrt{b})(a-\sqrt{b})=1$ 
		\item  $(3+b\sqrt{c})(3-b\sqrt{c})=1$
		\item  $(8+b\sqrt{c})(8-b\sqrt{c})=1$
		\item  $(a+b\sqrt{3})(a-b\sqrt{3})=1$
		\item  $(a+b\sqrt{5})(a-b\sqrt{5})=1$ ($b$ est différent de 0)
%		\item  $(a+b\sqrt{11})(a-b\sqrt{11})=1$
%		\item  $(a+b\sqrt{c})(a-b\sqrt{c})=4$ ($b$ et $c$ sont $\neq$ de 1 et 0)
%		\item  $(a+b\sqrt{c})(a-b\sqrt{c})=5$ 
%		\item  $(a+b\sqrt{c})(a-b\sqrt{c})=13$ 
	\end{enumerate}
\end{spacing}
\end{multicols}\vspace{-0.5em}
\end{exercise}
\begin{exercise} Mêmes consignes. Un peu plus dur !\vspace{-0.5em}
	\begin{multicols}{2} 
		\begin{spacing}{2}
			\begin{enumerate}[label={\arabic*)},leftmargin=*] 
				\item  $(a+b\sqrt{11})(a-b\sqrt{11})=1$
				\item  $(a+b\sqrt{c})(a-b\sqrt{c})=4$ ($b$ et $c$ sont $\neq$ de 1 et 0)
				\item  $(a+b\sqrt{c})(a-b\sqrt{c})=5$ 
				\item  $(a+b\sqrt{c})(a-b\sqrt{c})=13$ 
			\end{enumerate}
		\end{spacing}\vspace{-0.5em}
	\end{multicols}
\end{exercise}
\begin{exercise}[préliminaire]\label{chapter08:club:radicaux:identite} Montrer que pour tout $a$ et $b
\in\mathbb{R}$ on a $(a-b)^2=(a+b)^2-4ab$. 
\end{exercise}

\begin{eBox} Pour simplifier l'expression $\sqrt{52+16\sqrt{3}}$  nous chercherons deux  nombres (entiers ?) $a$ et $b$ tel que :
\setlength{\abovedisplayskip}{3pt}
\setlength{\belowdisplayskip}{3pt} 	
$$52+16\sqrt{3}=(\sqrt{a}+\sqrt{b})^2$$
\end{eBox}
\noindent \hfill\parbox{1\linewidth}{\centering
	$\begin{WithArrows}
%	52+16\sqrt{3}&=(\sqrt{a}+\sqrt{b})^2 \\
	52+ 2\sqrt{192}& = a + 2\sqrt{ab} + b 
\end{WithArrows}$ 
}
Il \textbf{suffit} de trouver $a$ et $b$ solutions du système \textbf{non linéaire} $ \begin{cases} a+b = 52 \\ ab=192 \end{cases}  $. Si \textbf{on ne peut pas deviner une solution}, on transformera le système à l'aide de l'identité de l'exercice \ref{chapter08:club:radicaux:identite} : 
\setlength{\abovedisplayskip}{3pt}
\setlength{\belowdisplayskip}{3pt}  
\[
(a-b)^2 =(a+b)^2-4ab = 52^2-4\times 192 = 1936  \qquad \text{donc} \qquad 	a-b   = 44
 \] 
  $a$ et $b$ sont aussi solutions de $\begin{cases} a+b = 52 \\ a-b=44 \end{cases}$, ce qui donne  $a=48$ et $b=4$.
\setlength{\abovedisplayskip}{3pt}
\setlength{\belowdisplayskip}{3pt}  
\[ 
52+16\sqrt{3}=(\sqrt{48}+\sqrt{4})^2=(4\sqrt{3}+2)^2 \qquad \text{donc} \qquad  \sqrt{52+16\sqrt{3}} = \abs{4\sqrt{3}+2} = 4\sqrt{3}+2 
\] 
\begin{exercise} \hfill 
	\begin{enumerate}[label={\arabic*)},leftmargin=*] 
	\item En suivant la démarche précédente trouver $a$ et $b$
	tel que $52+16\sqrt{3}=(\sqrt{a}+\sqrt{b})^2$.
	\item Simplifier au maximum l'expression  $\sqrt{52+16\sqrt{3}}$
\end{enumerate} 
\end{exercise}	 
\begin{exercise}[Entrainement] Déterminer les racines carrées de quelques radicaux parmi : 
	\vspace{-0.5em}
	\begin{multicols}{3} 
	\begin{enumerate}[label={$\Alph*=$},leftmargin=*] 
		\item 	$9+4\sqrt{5}$
		\item 	$25+4\sqrt{21}$
		\item 	$7-4\sqrt{3}$
		\item 	$8+2\sqrt{7}$
		\item 	$36-10\sqrt{11}$
		\item 	$36+16\sqrt{2}$
		\item 	$49+20\sqrt{6}$
		\item 	$29-12\sqrt{5}$
		\item 	$52+16\sqrt{3}$
		\item $64-24\sqrt{7}$
		\item $11+4\sqrt{6}$
		\item $8-2\sqrt{15}$
		\item $18-12\sqrt{2}$
		\item $55+30\sqrt{2}$
		\item $42-24\sqrt{3}$
	\end{enumerate}
	\end{multicols}	\vspace{-0.5em}
\end{exercise} 
\begin{exercise} Simplifier les expressions suivantes.	\vspace{-0.5em}
\begin{multicols}{3} 
	\begin{enumerate}[label={$\Alph*=$},leftmargin=*] 
	\item $\frac{1}{\sqrt{8+4\sqrt{3}}}$
	\item $\frac{1}{\sqrt{9+6\sqrt{2}}}$
	\item $\frac{1}{\sqrt{15-10\sqrt{2}}}$ 
\end{enumerate}\vspace{-0.5em}
\end{multicols}
\end{exercise}
\begin{exercise}\hfill \\
	 Calculer $\sqrt{1+\frac{1}{2}}\times 
	\sqrt{1+\frac{1}{3}}\times
	\sqrt{1+\frac{1}{4}}\times
	\sqrt{1+\frac{1}{5}}\times
	\ldots\times
	\sqrt{1+\frac{1}{100}}$.  
\end{exercise}
\begin{exercise}\hfill \\
	Calculer $(1+\sqrt{2})^{2022}(1-\sqrt{2})^{2023}$.
\end{exercise}
\begin{exercise}\hfill \\
	Calculer $(2\sqrt{5}+1)\left(\frac{1}{1+\sqrt{2}}+\frac{1}{\sqrt{2}+\sqrt{3}}+ \frac{1}{\sqrt{3}+\sqrt{4}}+\ldots+\frac{1}{\sqrt{99}+\sqrt{100}}\right)$
\end{exercise}
\begin{exercise} \hfill \\
	On pose $x=\frac{1}{2-\sqrt{3}}$. Déterminer la valeur de $x^2-4x+2$. Montre les calculs.
\end{exercise}
\begin{exercise} \hfill \\
	On pose $x=\sqrt{2}-1$. Déterminer la valeur de $x^3+2x^2+x+2$. Montre les calculs.
\end{exercise}
\begin{exercise} \hfill \\
	On pose $x=\sqrt{5}-2$. Déterminer la valeur de $(9+4\sqrt{5})x^2-(\sqrt{5}+2)x+4$. Montre les calculs.
\end{exercise}
\begin{exercise}\hfill \\ 
	Sachant que la somme de $(a-\sqrt{2})^2$ et $\sqrt{(b+1)^2}$ est égale à zéro, déterminer la valeur de $\frac{1}{b-a}$.
\end{exercise}
\begin{exercise}\hfill \\
	$m$ et $n$ sont deux nombres entier positifs. \\
	Rendre rationnel le dénominateur et simplifier l'expression $\frac{m+n-2\sqrt{mn}}{\sqrt{m}-\sqrt{n}}$.
\end{exercise}
\begin{exercise}\hfill \\
	Soit $x>1$. Rendre rationnel le dénominateur et simplifier l'expression $\frac{\sqrt{x+1}-\sqrt{x-1}}{\sqrt{x+1}+\sqrt{x-1}}$.
\end{exercise} 
\begin{exercise}\hfill \\ 
	On pose $x=4-\sqrt{2}$ et $y=4+\sqrt{2}$. Simplifier les expressions suivantes :\hfill 
	\begin{enumerate}[label={\arabic*)},leftmargin=*] 
		\item  $xy$ 
		\item $xy^2+x^2y$
		\item $\sqrt{x^2-2xy+y^2}$
	\end{enumerate}  
\end{exercise}
\begin{exercise} \hfill \\
	On pose $x=\sqrt{3}+\sqrt{2}$ et $y=\sqrt{3}-\sqrt{2}$. Simplifier les expressions suivantes :\hfill 
	\begin{enumerate}[label={\arabic*)},leftmargin=*] 
		\item  $x^2-xy+y^2$ 
		\item $x^3y+x^3y$ 
	\end{enumerate}  
\end{exercise}
\begin{exercise} \hfill \\
	Sachant que $a+\frac{1}{a}=1+\sqrt{10}$, simplifier $a^2+\frac{1}{a^2}$.
\end{exercise}
%----------------------------------------------------------------------------------------
%	
%----------------------------------------------------------------------------------------
\end{document}

 